\chapter{ToRCH Panel}

\section*{What Is This Test?}
The ToRCH panel is a group of serologic tests. It screens for antibodies to the most important congenital infections. These are infections transmitted from mother to fetus that can cause birth defects. The acronym TORCH (sometimes written ToRCH) stands for Toxoplasmosis, Other (syphilis, hepatitis B, HIV, varicella-zoster, parvovirus B19, and other agents), Rubella, Cytomegalovirus (CMV), and Herpes simplex virus (HSV). Congenital infection occurs transplacentally during pregnancy. It can cause miscarriage, stillbirth, intrauterine growth restriction, and a classic constellation of neonatal findings. These findings include microcephaly, intracranial calcifications, chorioretinitis, hepatosplenomegaly, jaundice, and sensorineural hearing loss. The panel measures IgM and IgG antibodies to each organism. IgM antibody in the newborn indicates congenital infection (maternal IgM does not cross the placenta). IgG antibody in the neonate is largely passively transferred maternal IgG. It is not by itself diagnostic. In the mother, the presence of IgG indicates past infection or vaccination (immunity). The presence of IgM suggests recent or primary infection. This carries the highest risk of fetal transmission. The panel is ordered when a pregnant woman has symptoms of, or exposure to, one of these infections. It is also ordered when fetal ultrasound shows concerning findings. It is also part of the evaluation of newborns with suspected congenital infection.

\section*{Alternative Names}
\begin{itemize}
  \item TORCH Screen
  \item ToRCH Serology
  \item TORCH Panel
  \item Congenital Infection Panel
  \item Congenital Toxoplasmosis, Other Agents, Rubella, CMV, HSV Screen
  \item TORCH Titers
  \item Maternal Infection Screen
\end{itemize}
\section*{Principle}
Each component is measured by serologic immunoassay. The most common methods are enzyme immunoassay (EIA) and chemiluminescent immunoassay (CLIA). These detect IgG and IgM antibodies specific to each organism. Immunoglobulin M is the first antibody produced after acute infection. It indicates recent or primary infection. Immunoglobulin G indicates past infection or immunity. In pregnancy, IgG may be passively transferred to the fetus. For toxoplasmosis and rubella, IgG avidity testing is used to time the infection. High-avidity IgG indicates an infection acquired at least 3--4 months earlier (low fetal risk). Low-avidity IgG indicates recent primary infection (higher fetal risk). For CMV, detection of the virus by PCR in amniotic fluid or neonatal blood is the definitive diagnostic test. This is because maternal IgM is frequently negative or nonspecific. For syphilis, the panel component uses nontreponemal (RPR) and treponemal tests. In the neonate, IgM detection indicates the infant's own antibody response. It therefore indicates congenital infection. A negative IgM does not fully exclude congenital infection. PCR-based testing (blood, urine, CSF, saliva) is preferred for CMV, HSV, and toxoplasmosis in the newborn. IgM assays are prone to false positives (rheumatoid factor, other infections, and nonspecific binding). Positive IgM results are typically confirmed by a second assay or by avidity testing.

\section*{History}
The term ``TORCH'' was coined in 1971 by Andr\'{e} Nahmias at a conference on congenital infections. It unified the group of infections that cause a similar syndrome of congenital disease. The individual infections were recognized earlier. Congenital syphilis was described in the 15th--16th centuries. The rubella virus was isolated in 1962. Its teratogenicity was established during the 1964--1965 rubella pandemic. This pandemic caused an estimated 20,000 congenital rubella syndrome cases in the United States. Congenital toxoplasmosis was characterized in the 1930s--1940s by Wolf, Cowen, and Paige. Cytomegalovirus was isolated in 1956. Its congenital role was recognized in the 1960s by Weller. Congenital herpes simplex infection was described in the 1930s by Batignani. The ``other'' agents (O) were added later. They include syphilis, hepatitis B, HIV, varicella-zoster, and parvovirus B19. The rubella vaccine (1969) and the dramatic reduction in congenital rubella syndrome validated the screening and prevention approach. Modern practice emphasizes organism-specific diagnosis (notably CMV PCR) over the single panel. This is because the panel's sensitivity and specificity for each agent vary considerably.

\section*{Why This Test Is Ordered}
\begin{itemize}
  \item \textbf{Maternal infection screening in pregnancy:} This is done when a pregnant woman develops a febrile illness, rash, lymphadenopathy, or other symptoms compatible with toxoplasmosis, rubella, CMV, or HSV. It is also done after known exposure to a child or adult with one of these infections
  \item \textbf{Routine prenatal screening:} Rubella IgG is universally assessed in early pregnancy to document immunity. Syphilis screening is standard prenatal care in all women. HIV and hepatitis B screening are universal in the United States
  \item \textbf{Fetal abnormalities on ultrasound:} Intracranial calcifications, ventriculomegaly, hydrops fetalis, growth restriction, or hepatosplenomegaly prompt a congenital infection workup
  \item \textbf{Evaluation of the newborn with suspected congenital infection:} Jaundice, hepatosplenomegaly, microcephaly, seizures, rash, chorioretinitis, or hearing loss prompt this evaluation
  \item \textbf{Evaluation of recurrent pregnancy loss or intrauterine fetal demise:} To identify a treatable or preventable infectious cause
  \item \textbf{Immunocompromised pregnant women:} Higher risk of primary infection and reactivation
\end{itemize}

\section*{Primary vs Secondary Test}
The ToRCH panel is a secondary (diagnostic) screening panel. It is not recommended as a routine screening tool for all pregnant women. The exceptions are the universal components: rubella IgG, syphilis, HIV, and hepatitis B. Routine screening for toxoplasmosis, CMV, and HSV is not currently recommended by most guidelines. The panel is ordered when a specific concern arises. Each positive result must be confirmed. It must be followed by organism-specific testing.

\section*{How This Test Is Performed}
A venous blood sample is collected in a red-top (serum separator) tube, typically 3--5 mL. It is tested by immunoassay for IgG and IgM to each of the panel organisms. Results are reported qualitatively (positive, negative, or equivocal). For rubella and some agents, results are also reported quantitatively as titers or index values. For toxoplasmosis and rubella, IgG avidity is measured when IgM is positive. This distinguishes recent from past infection. When a mother is found to be infected, fetal evaluation proceeds with targeted ultrasound. For CMV and toxoplasmosis, it also proceeds with amniocentesis and PCR of amniotic fluid. In the newborn, testing is performed on neonatal blood (IgM). Preferentially, PCR is performed on blood, urine, saliva, or CSF to detect the organism directly. Results are typically available within 1--3 days. PCR results are available within 24--72 hours.

\section*{What Preparation Is Needed}
No fasting or special preparation is required. The patient should inform the healthcare provider of gestational age. Also report prior rubella vaccination or known immunity. Report a history of prior HSV, CMV, or toxoplasmosis infections. Report occupation (childcare, farming, cat exposure) and dietary habits (undercooked meat). Report travel history. Report HIV or immunosuppressive conditions. Documentation of rubella immunization is important. Rubella IgG in a vaccinated woman indicates protective immunity. It eliminates the need for further testing.

\section*{Contraindications and Precautions}
\begin{itemize}
  \item There are no absolute contraindications to the blood draw. Important interpretive cautions:
  \item A positive maternal IgM does not always indicate recent primary infection (1). IgM can persist for months. It can be false-positive from rheumatoid factor, other infections, or autoantibodies. Confirmatory testing (IgG avidity, second assay, paired convalescent titers) is essential
  \item A negative maternal IgM does not exclude acute infection in immunocompromised patients (2)
  \item In the newborn, maternal IgG crosses the placenta and is present in the neonate (3). Only neonatal IgM indicates the infant's own infection. Even a negative IgM does not exclude congenital CMV or toxoplasmosis. So PCR is the preferred test
  \item CMV is the most common congenital infection (4). It is frequently asymptomatic in the mother. Maternal IgM is often negative or nonspecific. PCR is the diagnostic standard
  \item The panel is not a single unified test with uniform sensitivity (5). Each component differs. Negative results should not be overinterpreted
  \item Toxoplasmosis screening is recommended only in specific settings, not routinely (6). Primary prevention (avoiding undercooked meat and cat litter) is emphasized
  \item HSV IgM testing is particularly nonspecific (7). It is not recommended for routine screening. HSV serology is best used with type-specific assays
\end{itemize}
\section*{Risks}
Risks are those of routine venipuncture. They include mild pain, bruising, and a small hematoma at the collection site. They also include lightheadedness or vasovagal syncope. Extremely rarely, infection or nerve injury occurs. Amniocentesis for fetal PCR, when indicated, carries a small risk of pregnancy loss. It is performed only after counseling.

\section*{What Happens After the Test}
If all components are negative (no IgM; immune IgG where expected), the mother is counseled on primary prevention. The risk of infection is considered low. If a specific infection is identified, management is organism-specific. Acute toxoplasmosis is treated with spiramycin. For fetal infection, pyrimethamine-sulfadiazine with folinic acid is used. Primary rubella during pregnancy prompts counseling about the risk of congenital rubella syndrome. Termination is considered. CMV infection is evaluated with amniocentesis PCR and ultrasound. If severe, valganciclovir therapy is used. HSV is managed with suppressive acyclovir. Active genital lesions at delivery lead to cesarean section. Syphilis is treated with penicillin (the only fully effective regimen). A congenital infection in the neonate prompts a full evaluation (hearing screen, eye exam, imaging, PCR). Antiviral or antiparasitic therapy is given as indicated. The healthcare provider will coordinate with obstetrics, maternal-fetal medicine, infectious disease, and pediatrics.

\section*{Understanding Results}

\subsection*{Positive Maternal IgM with Low-Avidity IgG (Recent/Primary Infection)}
\begin{itemize}
  \item \textbf{Toxoplasmosis:} Risk of congenital infection varies with gestational age (highest in the third trimester). It is treated with spiramycin before fetal diagnosis
  \item \textbf{Rubella:} Highest risk is in the first trimester. Congenital rubella syndrome includes cardiac defects, cataracts, and sensorineural hearing loss
  \item \textbf{Cytomegalovirus (CMV):} Primary CMV carries the highest risk of symptomatic congenital infection. It is confirmed by amniotic fluid PCR
  \item \textbf{Herpes simplex virus (HSV):} Neonatal herpes is usually acquired at delivery rather than transplacentally. High-avidity IgM-persistent patterns are common. They are not diagnostic
\end{itemize}

\subsection*{Positive IgG Only (Past Infection or Immunity)}
\begin{itemize}
  \item \textbf{Toxoplasmosis IgG positive, IgM negative:} This indicates past infection. There is no increased fetal risk (unless immunocompromised)
  \item \textbf{Rubella IgG positive:} This indicates immunity (from infection or vaccination). There is no risk of congenital rubella syndrome
  \item \textbf{CMV IgG positive:} This indicates past or reactivated infection. There is low risk of symptomatic congenital infection
  \item \textbf{HSV type-specific IgG:} This indicates latent infection. The risk is neonatal transmission at delivery. This is managed by suppressive therapy and delivery planning
\end{itemize}

\subsection*{Negative Results}
\begin{itemize}
  \item \textbf{All IgM negative:} No evidence of acute infection with the tested agents at the time of sampling
  \item \textbf{No immunity (IgG negative to rubella):} The woman is rubella-susceptible. Vaccination postpartum is recommended. The live vaccine is contraindicated in pregnancy
  \item \textbf{CMV or toxoplasmosis negative serology:} Infection cannot be fully excluded. PCR-based testing is definitive when suspicion remains
\end{itemize}

\section*{Normal Results}
\begin{itemize}
  \item \textbf{Toxoplasma IgG and IgM:} Negative (non-immune) or IgG positive/IgM negative (past infection)
  \item \textbf{Rubella IgG:} Positive (immune, protective). Negative indicates susceptibility
  \item \textbf{CMV IgG and IgM:} Negative (non-immune) or IgG positive/IgM negative (past infection)
  \item \textbf{HSV-1 and HSV-2 type-specific IgG:} Negative or positive (latent infection). HSV IgM is not recommended for diagnosis
  \item \textbf{Other components (syphilis RPR, HIV, hepatitis B surface antigen):} Negative in an uninfected mother
  \item \textbf{Neonatal IgM:} Negative. A positive neonatal IgM indicates the infant's own antibody and congenital infection
\end{itemize}

\section*{Follow-up Tests}
\begin{itemize}
  \item \textbf{IgG avidity testing:} For toxoplasmosis and rubella, to time the infection when IgM is positive
  \item \textbf{PCR on amniotic fluid:} Definitive fetal diagnosis for CMV and toxoplasmosis
  \item \textbf{Targeted fetal ultrasound:} To assess for intracranial calcifications, hydrops, growth restriction, and structural anomalies
  \item \textbf{Neonatal PCR (blood, urine, saliva, CSF):} Preferred test for congenital CMV, HSV, and toxoplasmosis
  \item \textbf{Newborn hearing screen and ophthalmologic examination:} For chorioretinitis and sensorineural hearing loss
  \item \textbf{Neuroimaging (cranial ultrasound or MRI):} For intracranial calcifications and periventricular changes
  \item \textbf{Repeat serology (paired titers):} 2--4 weeks later to document rising titers
  \item \textbf{RPR and confirmatory treponemal tests:} For syphilis. Treat with penicillin
  \item \textbf{Rubella vaccination:} Postpartum for rubella-susceptible women (live vaccine contraindicated in pregnancy)
\end{itemize}

\section*{References}
\begin{enumerate}
  \item Nahmias AJ, Walls KW, Stewart JA, Herrmann KL, Flynn WJ. The ToRCH complex-perinatal infections associated with toxoplasma and rubella, cytomegalo- and herpes simplex viruses. Pediatr Res. 1971;5(8):405--406.
  \item ACOG Practice Bulletin No. 151: Cytomegalovirus, parvovirus B19, varicella zoster, and toxoplasmosis in pregnancy. Obstet Gynecol. 2015;125(6):1510--1525.
  \item ACOG Practice Bulletin No. 207: Thrombocytopenia in pregnancy. Obstet Gynecol. 2019;133(3):e181--e193.
  \item Montoya JG, Liesenfeld O. Toxoplasmosis. Lancet. 2004;363(9425):1965--1976.
  \item Dontigny L, Arsenault MY, Martel MJ, et al. Rubella in pregnancy. J Obstet Gynaecol Can. 2008;30(2):152--158.
  \item Kimberlin DW. Neonatal herpes simplex infection. Clin Microbiol Rev. 2004;17(1):1--13.
  \item Rawlinson WD, Boppana SB, Fowler KB, et al. Congenital cytomegalovirus infection in pregnancy and the neonate: consensus recommendations for prevention, diagnosis, and therapy. Lancet Infect Dis. 2017;17(6):e177--e188.
  \item Wallach J. Interpretation of Diagnostic Tests. 11th ed. Wolters Kluwer; 2022.
\end{enumerate}

\clearpage
\section*{Regulatory Status and Authorization Date}
This chapter describes a test category, not a specific commercial product. FDA, CDSCO, EU IVDR, and MHRA approval or registration dates therefore cannot be assigned without the assay manufacturer, product name, instrument, and intended use. For laboratory-developed tests, an individual agency approval date may not exist. See the Preface for the interpretation of regulatory status.

\clearpage
